% ================= Topic 3 (VI) -- part d: problems 41--50 (nâng cao) =================
\begin{probbox}{41}{\stars{1}}
Trên một lưới điểm vuông, bạn Ruby đứng ở điểm $A$. Ruby tung một đồng xu cân đối $6$ lần; mỗi lần ra mặt sấp thì Ruby đi
\emph{lên} một điểm, ra mặt ngửa thì Ruby đi \emph{sang phải} một điểm (như bài 16). Các điểm $M$, $N$, $B$ được đánh dấu
trên hình.
\begin{figC}\begin{scope}[scale=0.85]
  \foreach \x in {0,...,3}{\foreach \y in {0,...,3}{\fill[Navy] (\x,\y) circle (2.4pt);}}
  \fill[tmOrange] (0,0) circle (3.4pt);\node[below left] at (0,0) {$A$};
  \fill[tmGreen!80!black] (1,1) circle (3.4pt);\node[above left] at (1,1) {$M$};
  \fill[AnsRed] (2,1) circle (3.4pt);\node[below right] at (2,1) {$N$};
  \fill[tmPurple] (3,3) circle (3.4pt);\node[above right] at (3,3) {$B$};
\end{scope}\end{figC}
\begin{enumerate}[label=\alph*),leftmargin=1.6em,itemsep=1pt,topsep=2pt]
\item Tính xác suất để sau $6$ lần tung, Ruby dừng ở điểm $B$.
\item Tính xác suất để Ruby đi qua điểm $M$ \emph{và} dừng ở $B$.
\item Tính xác suất để Ruby dừng ở $B$ mà \emph{không} đi qua điểm $N$.
\end{enumerate}
\end{probbox}
\needspace{8\baselineskip}
\probsol{
\textbf{Phân tích đề.} Mỗi dãy $6$ lần tung là một dãy gồm $6$ kí hiệu $\uparrow$ (lên) hoặc $\rightarrow$ (sang phải).
Có $2^6=2\times2\times2\times2\times2\times2=64$ dãy, đồng khả năng. Mỗi dãy là một \emph{đường đi} trên lưới.
Ta đánh dấu vị trí bằng cặp (số bước sang phải, số bước lên): $A=(0,0)$, $M=(1,1)$, $N=(2,1)$, $B=(3,3)$.

\step{1}{Đếm số đường đi từ $A$ tới mỗi điểm (câu a)}
\begin{recallbox}Một điểm chỉ có thể tới được từ điểm ngay bên trái hoặc ngay bên dưới nó. Vì vậy:
số đường tới một điểm $=$ số đường tới điểm bên trái $+$ số đường tới điểm bên dưới (quy tắc cộng).\end{recallbox}
Ghi số đường đi vào từng điểm, bắt đầu từ $A$ (ghi $1$); hàng dưới cùng và cột trái cùng đều là $1$:
\begin{figC}\begin{scope}[scale=0.85]
  \foreach \x/\y/\n in {0/0/1,1/0/1,2/0/1,3/0/1,0/1/1,1/1/2,2/1/3,3/1/4,0/2/1,1/2/3,2/2/6,3/2/10,0/3/1,1/3/4,2/3/10,3/3/20}{
    \fill[tmBlue!12] (\x,\y) circle (9pt);\node[font=\small\bfseries,text=Navy] at (\x,\y) {\n};}
  \draw[tmPurple,thick] (3,3) circle (10pt);\node[right=9pt,text=tmPurple] at (3,3) {$B$};
\end{scope}\end{figC}
Có $20$ đường đi từ $A$ tới $B$. (Cách khác: đi tới $B$ cần đúng $3$ bước lên trong $6$ bước; số cách chọn $3$ lượt
trong $6$ lượt là $\dfrac{6\times5\times4}{3\times2\times1}=20$.)
\[P(a)=\frac{20}{64}=\frac{5}{16}.\]

\step{2}{Câu b: đi qua $M$ rồi tới $B$ -- quy tắc nhân}
Một đường qua $M$ gồm hai chặng: chặng $A\to M$ ($1$ bước lên, $1$ bước phải) và chặng $M\to B$ ($2$ bước lên, $2$ bước phải).
Chặng $A\to M$ có $2$ cách ($\uparrow\rightarrow$ hoặc $\rightarrow\uparrow$); chặng $M\to B$ có $6$ cách (giống bài 16:
chọn $2$ lượt lên trong $4$ lượt). Mỗi cách đi chặng đầu ghép với mỗi cách đi chặng sau cho một đường khác nhau:
$2\times6=12$ đường.
\[P(b)=\frac{12}{64}=\frac{3}{16}.\]

\step{3}{Câu c: không đi qua $N$ -- dùng phần bù}
Số đường từ $A$ qua $N$ tới $B$: chặng $A\to N$ ($2$ bước phải, $1$ bước lên) có $3$ cách; chặng $N\to B$ ($1$ bước phải,
$2$ bước lên) có $3$ cách; tổng $3\times3=9$ đường. Vậy số đường từ $A$ tới $B$ \emph{không} qua $N$ là $20-9=11$.
\[P(c)=\frac{11}{64}.\]
\begin{proofbox}\emph{Kiểm tra bằng cách ``chặn'' điểm $N$.} Ghi lại số đường đi nhưng đặt số $0$ tại $N$ (coi như không được đi qua).
Hàng $1$: $1,\ 2,\ \mathbf{0},\ 0+1=1$. Hàng $2$: $1,\ 3,\ 3+0=3,\ 3+1=4$. Hàng $3$: $1,\ 4,\ 4+3=7,\ 7+4=11$.
Tại $B$ được $11$, khớp với phần bù.\end{proofbox}
\begin{warnbox}Câu b và c hỏi về xác suất trong \emph{toàn bộ} $64$ dãy tung, không phải trong $20$ đường tới $B$. Nếu chia cho $20$
là trả lời một câu hỏi khác (xác suất có điều kiện, chưa học).\end{warnbox}
\verify{Ở câu a, các điểm cuối có thể là $(6,0),(5,1),\dots,(0,6)$ với số đường $1,6,15,20,15,6,1$; tổng $=64$. Điểm $B=(3,3)$ ứng với số $20$.}
\ans{a) $\dfrac5{16}$;\quad b) $\dfrac3{16}$;\quad c) $\dfrac{11}{64}$}
}

\begin{probbox}{42}{\stars{2}}
Ta xếp các lưới ô vuông bằng que tăm như ở bài 17. Một lưới ``$m\times n$'' có $m$ hàng, mỗi hàng $n$ ô vuông.
\begin{enumerate}[label=\alph*),leftmargin=1.6em,itemsep=1pt,topsep=2pt]
\item Lập công thức tính số que tăm cần dùng cho lưới $m\times n$.
\item Một lưới hình vuông $k\times k$ dùng hết đúng $760$ que tăm. Tìm $k$.
\item Bạn Sam có đúng $2026$ que tăm và muốn xếp một lưới $m\times n$ (với $1\le m\le n$) dùng \emph{hết} số que này.
Tìm tất cả các lưới Sam có thể xếp.
\end{enumerate}
\end{probbox}
\probsol{
\textbf{Phân tích đề.} Giống bài 17, ta đếm riêng que nằm ngang và que thẳng đứng. Câu c dẫn tới một phương trình với hai ẩn
là số tự nhiên; mẹo ``nhân $2$ rồi cộng $1$'' giúp biến nó thành bài toán tìm ước.

\step{1}{Câu a: công thức}
Lưới $m\times n$ có $m+1$ đường kẻ ngang, mỗi đường $n$ que: $(m+1)\times n$ que nằm ngang. Có $n+1$ đường kẻ dọc, mỗi đường
$m$ que: $(n+1)\times m$ que thẳng đứng. Tổng:
\[T=(m+1)n+(n+1)m=mn+n+mn+m=2mn+m+n.\]
Thử lại với lưới $1\times10$: $2\times1\times10+1+10=31$. Đúng như bài 17.

\step{2}{Câu b: lưới vuông $k\times k$}
Thay $m=n=k$: $T=2k^2+2k=2k(k+1)$. Ta cần $2k(k+1)=760$, tức $k(k+1)=380$.
$k(k+1)$ là tích hai số tự nhiên liên tiếp. Vì $19\times20=380$, nên $k=19$. (Khi $k$ tăng thì $k(k+1)$ tăng, nên chỉ có một
giá trị: $18\times19=342<380<420=20\times21$.)

\step{3}{Câu c: lập phương trình}
Cần tìm các số tự nhiên $1\le m\le n$ với $2mn+m+n=2026$.
\begin{tipbox}Nhân hai vế với $2$ rồi cộng $1$: $4mn+2m+2n+1=4053$. Vế trái phân tích được thành tích:
$(2m+1)(2n+1)=4mn+2m+2n+1$ (nhân phân phối để kiểm tra). Vậy
\[(2m+1)(2n+1)=4053.\]\end{tipbox}

\step{4}{Phân tích $4053$ ra thừa số nguyên tố}
$4053$ có tổng chữ số $4+0+5+3=12$ chia hết cho $3$: $4053=3\times1351$. Tiếp theo $1351=7\times193$.
Số $193$ là số nguyên tố: nó không chia hết cho $2,3,5,7,11,13$, và $13\times13=169<193<196=14\times14$, nên không cần thử thêm.
\begin{recallbox}Để kiểm tra $N$ có là số nguyên tố không, chỉ cần thử chia cho các số nguyên tố $p$ với $p\times p\le N$.\end{recallbox}
Vậy $4053=3\times7\times193$. Các ước của $4053$ là: $1,\ 3,\ 7,\ 21,\ 193,\ 579,\ 1351,\ 4053$.

\step{5}{Tìm các cặp $(m,n)$}
Vì $m\ge1$ nên $2m+1\ge3$; vì $m\le n$ nên $2m+1\le2n+1$. Các cách viết $4053$ thành tích hai ước thỏa mãn:
\begin{center}\begin{tabular}{|c|c|c|c|c|}\hline
$2m+1$ & $2n+1$ & $m$ & $n$ & Thử lại $2mn+m+n$\\\hline
$3$ & $1351$ & $1$ & $675$ & $1350+1+675=2026$\\\hline
$7$ & $579$ & $3$ & $289$ & $1734+3+289=2026$\\\hline
$21$ & $193$ & $10$ & $96$ & $1920+10+96=2026$\\\hline
\end{tabular}\end{center}
(Cặp $1\times4053$ bị loại vì cho $m=0$.)
\begin{warnbox}Thử từng giá trị $m=1,2,3,\dots$ cũng làm được nhưng rất dài, dễ sót. Mẹo nhân $2$ cộng $1$ biến bài toán thành
``tìm ước'' và đảm bảo \emph{không bỏ sót} nghiệm nào.\end{warnbox}
\verify{Ở các giá trị $m$ khác, $2m+1$ không phải là ước của $4053$ nên $n$ không là số tự nhiên. Ví dụ $m=2$: $5(2n+1)=4053$ không
được vì $4053$ không chia hết cho $5$.}
\ans{a) $2mn+m+n$;\quad b) $k=19$;\quad c) ba lưới: $1\times675$, $3\times289$, $10\times96$}
}

\begin{probbox}{43}{\stars{1}}
Gieo cùng lúc ba con xúc xắc cân đối có màu đỏ, xanh và vàng (mỗi con ghi số từ $1$ đến $6$). Tính xác suất của các biến cố:
\begin{enumerate}[label=\alph*),leftmargin=1.6em,itemsep=1pt,topsep=2pt]
\item ba số xuất hiện đôi một khác nhau;
\item số trên con đỏ nhỏ hơn số trên con xanh, và số trên con xanh nhỏ hơn số trên con vàng;
\item có ít nhất một con xúc xắc ra số $6$;
\item số lớn nhất trong ba số xuất hiện bằng đúng $4$.
\end{enumerate}
\end{probbox}
\probsol{
\textbf{Phân tích đề.} Mỗi kết quả là bộ ba (đỏ, xanh, vàng). Theo quy tắc nhân có $6\times6\times6=216$ kết quả đồng khả năng.

\step{1}{Câu a: ba số khác nhau}
Con đỏ: $6$ cách; con xanh khác con đỏ: $5$ cách; con vàng khác hai con kia: $4$ cách. Có $6\times5\times4=120$ kết quả.
\[P(a)=\frac{120}{216}=\frac{5}{9}.\]

\step{2}{Câu b: đỏ $<$ xanh $<$ vàng -- dùng tính đối xứng}
Xét một bộ ba số khác nhau, chẳng hạn $\{2,3,5\}$. Có $3\times2\times1=6$ cách xếp ba số này cho ba con xúc xắc
($235$, $253$, $325$, $352$, $523$, $532$), và \emph{chỉ đúng một} cách là tăng dần ($2<3<5$).
Như vậy $120$ kết quả ở câu a chia thành các nhóm $6$, mỗi nhóm có $1$ kết quả thuận lợi: số kết quả thuận lợi là $120:6=20$.
\[P(b)=\frac{20}{216}=\frac{5}{54}.\]
\begin{recallbox}Số cách xếp $3$ vật khác nhau vào $3$ chỗ là $3\times2\times1=6$.\end{recallbox}

\step{3}{Câu c: ít nhất một số $6$ -- phần bù}
Biến cố đối là ``không con nào ra $6$'': mỗi con chỉ có $5$ kết quả ($1$ đến $5$), tổng $5\times5\times5=125$.
Vậy số kết quả có ít nhất một số $6$ là $216-125=91$.
\[P(c)=\frac{91}{216}.\]
\begin{warnbox}Không được tính ``mỗi con có xác suất $\frac16$ ra số $6$, ba con là $\frac36=\frac12$''. Cách cộng đó đếm trùng các kết quả
có hai hoặc ba số $6$ (như $(6,6,1)$ bị đếm hai lần).\end{warnbox}

\step{4}{Câu d: số lớn nhất bằng đúng $4$}
``Số lớn nhất $\le4$'' nghĩa là cả ba số đều thuộc $\{1,2,3,4\}$: $4\times4\times4=64$ kết quả.
``Số lớn nhất $\le3$'': $3\times3\times3=27$ kết quả. Số lớn nhất bằng đúng $4$ khi nó $\le4$ nhưng không $\le3$:
$64-27=37$ kết quả.
\[P(d)=\frac{37}{216}.\]
\begin{proofbox}Mọi kết quả có số lớn nhất $\le3$ đều nằm trong nhóm có số lớn nhất $\le4$. Bớt đi nhóm nhỏ còn lại đúng các kết
quả mà số lớn nhất không vượt quá $4$ nhưng vượt quá $3$, tức là bằng $4$.\end{proofbox}
\verify{Cách nghĩ của câu d áp dụng cho câu c: số lớn nhất bằng $6$ có $6^3-5^3=216-125=91$ kết quả, khớp. Tổng số kết quả theo
số lớn nhất $1,2,\dots,6$: $1+7+19+37+61+91=216$.}
\ans{a) $\dfrac59$;\quad b) $\dfrac5{54}$;\quad c) $\dfrac{91}{216}$;\quad d) $\dfrac{37}{216}$}
}

\begin{probbox}{44}{\stars{2}}
Người ta làm các con xúc xắc đặc biệt có $6$ mặt, mỗi mặt ghi một số trong các số $1,2,3,4,5,6$ (được phép ghi trùng), sao cho
giống bài 19: xác suất gieo được số chia hết cho $3$ là $\frac12$ và xác suất gieo được số chẵn là $\frac13$.
Hai con xúc xắc được coi là như nhau nếu chúng có cùng bộ sáu số (không kể thứ tự các mặt).
\begin{enumerate}[label=\alph*),leftmargin=1.6em,itemsep=1pt,topsep=2pt]
\item Tổng sáu số trên một con xúc xắc như vậy lớn nhất là bao nhiêu? Nhỏ nhất là bao nhiêu?
\item Có tất cả bao nhiêu con xúc xắc khác nhau như vậy?
\end{enumerate}
\end{probbox}
\probsol{
\textbf{Phân tích đề.} Như bài 19, điều kiện nghĩa là: \textbf{đúng $3$ mặt chia hết cho $3$} và \textbf{đúng $2$ mặt chẵn}.
Mấu chốt là số $6$ thuộc cả hai nhóm, nên ta phân loại các số $1$--$6$ thành bốn loại:
\begin{center}\begin{tabular}{|c|c|c|}\hline
Loại & Các số & Tính chất\\\hline
I & $6$ & chẵn và chia hết cho $3$\\\hline
II & $3$ & lẻ, chia hết cho $3$\\\hline
III & $2,\ 4$ & chẵn, không chia hết cho $3$\\\hline
IV & $1,\ 5$ & lẻ, không chia hết cho $3$\\\hline
\end{tabular}\end{center}

\step{1}{Biểu diễn số mặt mỗi loại theo số mặt ghi số $6$}
Gọi $a$ là số mặt ghi số $6$. Khi đó:
\begin{itemize}[leftmargin=1.2em,itemsep=1pt,topsep=2pt]
\item số mặt chia hết cho $3$ là (loại I) $+$ (loại II) $=3$, nên loại II có $3-a$ mặt;
\item số mặt chẵn là (loại I) $+$ (loại III) $=2$, nên loại III có $2-a$ mặt;
\item loại IV có $6-a-(3-a)-(2-a)=1+a$ mặt.
\end{itemize}
Số mặt không thể âm nên $2-a\ge0$, tức $a\in\{0,1,2\}$.
\begin{center}\begin{tabular}{|c|c|c|c|c|}\hline
$a$ & loại I (số $6$) & loại II (số $3$) & loại III ($2$ hoặc $4$) & loại IV ($1$ hoặc $5$)\\\hline
$0$ & $0$ & $3$ & $2$ & $1$\\\hline
$1$ & $1$ & $2$ & $1$ & $2$\\\hline
$2$ & $2$ & $1$ & $0$ & $3$\\\hline
\end{tabular}\end{center}

\step{2}{Câu a: tổng lớn nhất và nhỏ nhất}
Muốn tổng lớn nhất, ở loại III ta chọn toàn số $4$, ở loại IV chọn toàn số $5$. Tổng là
\[6a+3(3-a)+4(2-a)+5(1+a)=22+4a,\]
lớn nhất khi $a=2$: tổng $30$ (con xúc xắc $3,5,5,5,6,6$).
Muốn tổng nhỏ nhất, chọn toàn số $2$ ở loại III và toàn số $1$ ở loại IV. Tổng là
\[6a+3(3-a)+2(2-a)+1\cdot(1+a)=14+2a,\]
nhỏ nhất khi $a=0$: tổng $14$ (con xúc xắc $1,2,2,3,3,3$).
\begin{recallbox}Thu gọn: $6a-3a-4a+5a=4a$ và $9+8+5=22$; tương tự $6a-3a-2a+a=2a$ và $9+4+1=14$.\end{recallbox}

\step{3}{Câu b: đếm số con xúc xắc}
Khi đã biết $a$, ta chỉ còn phải quyết định: trong các mặt loại III có bao nhiêu mặt là số $4$ (còn lại là số $2$), và trong các mặt loại IV
có bao nhiêu mặt là số $5$ (còn lại là số $1$).
\begin{itemize}[leftmargin=1.2em,itemsep=1pt,topsep=2pt]
\item $a=0$: loại III có $2$ mặt, số mặt ghi $4$ có thể là $0,1,2$: $3$ cách; loại IV có $1$ mặt: $2$ cách. Được $3\times2=6$ con.
\item $a=1$: loại III có $1$ mặt: $2$ cách; loại IV có $2$ mặt: $3$ cách. Được $2\times3=6$ con.
\item $a=2$: loại III có $0$ mặt: $1$ cách; loại IV có $3$ mặt (số mặt ghi $5$ là $0,1,2,3$): $4$ cách. Được $1\times4=4$ con.
\end{itemize}
Tổng: $6+6+4=16$ con xúc xắc.
\begin{warnbox}Vì không kể thứ tự các mặt, hai mặt loại III là ``$2,4$'' hay ``$4,2$'' chỉ là \emph{một} trường hợp. Đếm số
\emph{lượng} số $4$ (là $0$, $1$ hay $2$) thay vì xếp chỗ sẽ tránh đếm trùng.\end{warnbox}
\verify{Hai con trong danh sách: với $a=1$, chọn một số $4$ và hai số $1$: con xúc xắc $1,1,3,3,4,6$ có $3$ mặt chia hết cho $3$
($3,3,6$) và $2$ mặt chẵn ($4,6$). Đúng. Con $1,2,2,3,3,3$ (tổng $14$) cũng có $3$ mặt chia hết cho $3$ và $2$ mặt chẵn.}
\ans{a) Lớn nhất $30$, nhỏ nhất $14$;\quad b) $16$ con xúc xắc}
}

\begin{probbox}{45}{\stars{1}}
Ba ngọn đèn nhấp nháy: đèn thứ nhất cứ $4$ giây nháy một lần, đèn thứ hai cứ $6$ giây, đèn thứ ba cứ $10$ giây. Lúc bắt đầu
($0$ giây) cả ba đèn cùng nháy. Xét $10$ phút tiếp theo (từ giây thứ $1$ đến hết giây thứ $600$).
\begin{enumerate}[label=\alph*),leftmargin=1.6em,itemsep=1pt,topsep=2pt]
\item Có bao nhiêu lần cả ba đèn cùng nháy?
\item Có bao nhiêu thời điểm có \emph{đúng hai} đèn cùng nháy?
\item Có bao nhiêu thời điểm có ít nhất một đèn nháy?
\end{enumerate}
\end{probbox}
\probsol{
\textbf{Phân tích đề.} Đèn thứ nhất nháy ở các giây là bội của $4$; đèn thứ hai ở các bội của $6$; đèn thứ ba ở các bội của $10$
(ý tưởng giống bài 20: các thời điểm ``trùng nhau được tìm bằng bội chung). Ta đếm các bội trong khoảng từ $1$ đến $600$.

\step{1}{Đếm số lần mỗi đèn nháy}
\begin{recallbox}Số bội của $d$ trong các số từ $1$ đến $N$ là phần nguyên của $N:d$ (thương trong phép chia có dư).\end{recallbox}
Đèn 1: $600:4=150$ lần. Đèn 2: $600:6=100$ lần. Đèn 3: $600:10=60$ lần.

\step{2}{Hai đèn và ba đèn cùng nháy -- dùng BCNN}
Hai đèn cùng nháy khi số giây là bội chung của hai chu kì, tức là bội của BCNN.
\begin{itemize}[leftmargin=1.2em,itemsep=1pt,topsep=2pt]
\item $\text{BCNN}(4,6)=12$: $600:12=50$ lần đèn 1 và đèn 2 cùng nháy.
\item $\text{BCNN}(4,10)=20$: $600:20=30$ lần.
\item $\text{BCNN}(6,10)=30$: $600:30=20$ lần.
\item $\text{BCNN}(4,6,10)=60$: $600:60=10$ lần cả ba đèn cùng nháy. Đó là đáp số câu a.
\end{itemize}
\begin{warnbox}$\text{BCNN}(4,6)$ là $12$, không phải $4\times6=24$ (vì $4$ và $6$ có ước chung $2$). Tương tự
$\text{BCNN}(4,6,10)=60$ chứ không phải $240$.\end{warnbox}

\step{3}{Câu b: đúng hai đèn}
Trong $50$ lần đèn 1 và đèn 2 cùng nháy có $10$ lần đèn 3 cũng nháy (ba đèn), nên chỉ có hai đèn 1, 2: $50-10=40$ lần.
Tương tự: chỉ đèn 1 và đèn 3: $30-10=20$; chỉ đèn 2 và đèn 3: $20-10=10$. Tổng số thời điểm đúng hai đèn:
$40+20+10=70$.

\step{4}{Câu c: vẽ sơ đồ Ven và điền số từ trong ra ngoài}
Chỉ đèn 1: $150-40-20-10=80$; chỉ đèn 2: $100-40-10-10=40$; chỉ đèn 3: $60-20-10-10=20$.
\begin{figC}
  \draw[tmBlue,thick] (0,0) circle (1.3);\draw[tmOrange,thick] (1.6,0) circle (1.3);\draw[tmGreen!80!black,thick] (0.8,-1.3) circle (1.3);
  \node[text=tmBlue,anchor=south east,inner sep=1pt] at (-0.95,0.95) {đèn 1};\node[text=tmOrange!80!black,anchor=south west,inner sep=1pt] at (2.55,0.95) {đèn 2};
  \node[text=tmGreen!70!black,anchor=north,inner sep=2pt] at (0.8,-2.62) {đèn 3};
  \node at (-0.45,0.25) {$80$};\node at (2.05,0.25) {$40$};\node at (0.8,-1.95) {$20$};
  \node at (0.8,0.5) {$40$};\node at (0.05,-0.95) {$20$};\node at (1.55,-0.95) {$10$};
  \node[text=AnsRed,font=\bfseries\small] at (0.8,-0.45) {$10$};
\end{figC}
Số thời điểm có ít nhất một đèn nháy là tổng bảy số trong sơ đồ:
\[80+40+20+40+20+10+10=220.\]
\begin{recallbox}\emph{Nguyên lí bù trừ} cho ba nhóm: số phần tử thuộc ít nhất một nhóm $=$ (tổng số phần tử của ba nhóm) $-$ (tổng số phần tử chung của từng cặp hai nhóm) $+$ (số phần tử chung của cả ba nhóm). Phần chung của cả ba nhóm bị cộng $3$ lần, bị trừ $3$ lần, nên phải cộng lại $1$ lần.\end{recallbox}
\verify{Theo nguyên lí bù trừ: $150+100+60-50-30-20+10=220$. Khớp với sơ đồ Ven.}
\ans{a) $10$ lần;\quad b) $70$ thời điểm;\quad c) $220$ thời điểm}
}


\begin{probbox}{46}{\stars{2}}
\begin{enumerate}[label=\alph*),leftmargin=1.6em,itemsep=2pt,topsep=2pt]
\item Một hình chóp tứ giác đều có $5$ mặt: một mặt đáy hình vuông và bốn mặt bên hình tam giác. Cần ít nhất bao nhiêu màu để sơn
$5$ mặt (mỗi mặt một màu) sao cho hai mặt có chung một cạnh luôn khác màu?
\item Có $6$ màu khác nhau. Ta sơn sáu mặt của một hình lập phương, mỗi mặt một màu, dùng đủ cả $6$ màu. Hai hình lập phương được
coi là giống nhau nếu có thể xoay hình này để được hình kia (màu ở các vị trí trùng khớp). Có bao nhiêu hình lập phương khác nhau?
\end{enumerate}
\begin{figC}
  \begin{scope}
    \coordinate (A) at (0,0);\coordinate (B) at (2,0);\coordinate (C) at (2.7,0.6);\coordinate (D) at (0.7,0.6);\coordinate (S) at (1.35,2.1);
    \fill[tmBlue!30] (A)--(B)--(S)--cycle;\fill[tmOrange!35] (B)--(C)--(S)--cycle;
    \draw[Ink,thick,dashed] (A)--(D)--(C);\draw[Ink,thick,dashed] (D)--(S);
    \draw[Ink,thick] (A)--(B)--(C)--(S)--cycle;\draw[Ink,thick] (B)--(S);
    \node[font=\scriptsize] at (1.35,-0.3) {hình chóp tứ giác};
  \end{scope}
  \begin{scope}[shift={(4.2,0)}]
    \fill[tmGreen!35,draw=Ink,thick] (0,0)--(1.4,0)--(1.4,1.4)--(0,1.4)--cycle;
    \fill[tmPurple!30,draw=Ink,thick] (0,1.4)--(1.4,1.4)--(2.0,1.95)--(0.6,1.95)--cycle;
    \fill[tmTeal!30,draw=Ink,thick] (1.4,0)--(2.0,0.55)--(2.0,1.95)--(1.4,1.4)--cycle;
    \node[font=\scriptsize] at (1,-0.3) {hình lập phương $6$ màu};
  \end{scope}
\end{figC}
\end{probbox}
\probsol{
\textbf{Phân tích đề.} Câu a giống bài 21: tìm một cách tô và chứng minh không thể ít màu hơn. Câu b là bài toán đếm ``không kể
cách xoay'': phải cẩn thận vì một hình lập phương có thể đặt theo nhiều tư thế.

\step{1}{Câu a: cần ít nhất $3$ màu}
\begin{proofbox}Mặt đáy có chung cạnh với cả bốn mặt bên, nên màu của đáy khác màu của mọi mặt bên. Hai mặt bên kề nhau cũng có chung
một cạnh (cạnh bên đi lên đỉnh chóp), nên cũng phải khác màu. Lấy mặt đáy và hai mặt bên kề nhau: ba mặt này đôi một chung cạnh,
nên cần ít nhất $3$ màu.\end{proofbox}
\emph{Ba màu là đủ:} sơn đáy màu $1$; bốn mặt bên đi vòng quanh sơn xen kẽ màu $2$, màu $3$, màu $2$, màu $3$. Mặt bên thứ tư (màu $3$)
kề với mặt bên thứ nhất (màu $2$): khác màu. Vậy câu a cần ít nhất $3$ màu.
\begin{tipbox}Nếu đáy là ngũ giác (hình chóp có $5$ mặt bên) thì $5$ mặt bên xếp thành vòng lẻ, không thể tô xen kẽ hai màu,
nên cần $4$ màu. Số mặt bên chẵn hay lẻ quyết định đáp số!\end{tipbox}

\step{2}{Câu b: đặt hình lập phương về một ``tư thế chuẩn''}
Gọi các màu là $1,2,3,4,5,6$. Với mỗi hình lập phương, ta xoay nó theo hai bước:
\begin{itemize}[leftmargin=1.2em,itemsep=1pt,topsep=2pt]
\item xoay sao cho mặt màu $1$ nằm \emph{dưới đáy};
\item giữ nguyên đáy, xoay quanh trục thẳng đứng sao cho, trong bốn mặt bên, mặt có \emph{số màu nhỏ nhất} nằm ở phía trước.
\end{itemize}
Sau khi đặt như vậy, không còn phép xoay nào giữ nguyên cả mặt đáy lẫn mặt trước. Do đó hai hình lập phương là giống nhau khi và
chỉ khi ở tư thế chuẩn chúng có cùng màu ở mọi vị trí.

\step{3}{Đếm số tư thế chuẩn}
\begin{itemize}[leftmargin=1.2em,itemsep=1pt,topsep=2pt]
\item Mặt trên (đối diện màu $1$): chọn một trong $5$ màu còn lại: $5$ cách.
\item Bốn màu còn lại nằm ở bốn mặt bên; màu nhỏ nhất trong bốn màu đó đã được đặt ở mặt trước (không phải chọn).
\item Ba màu còn lại xếp vào ba mặt: trái, sau, phải: $3\times2\times1=6$ cách.
\end{itemize}
Theo quy tắc nhân: $5\times6=30$ hình lập phương khác nhau.

\step{4}{Kiểm tra bằng cách thứ hai}
Nếu \emph{không} cho phép xoay (sáu vị trí trên, dưới, trước, sau, trái, phải là cố định), số cách sơn là
$6\times5\times4\times3\times2\times1=720$. Mỗi hình lập phương có thể đặt theo $6\times4=24$ tư thế ($6$ cách chọn mặt nằm dưới
đáy, rồi $4$ cách xoay quanh trục thẳng đứng), và với $6$ màu khác nhau thì $24$ tư thế này cho $24$ cách sơn cố định khác nhau.
Vậy số hình lập phương khác nhau là $720:24=30$, khớp với cách đếm ở bước 3.
\begin{warnbox}Trả lời $720$ là quên rằng các hình xoay được thành nhau thì giống nhau. Trả lời $6\times5=30$ ``cho có'' mà không giải thích
cũng chưa đủ: cần lập luận vì sao mỗi hình lập phương được đếm đúng một lần.\end{warnbox}
\ans{a) $3$ màu;\quad b) $30$ hình lập phương}
}

\begin{probbox}{47}{\stars{2}}
Xét các số nguyên $1,2,3,\dots,20$.
\begin{enumerate}[label=\alph*),leftmargin=1.6em,itemsep=1pt,topsep=2pt]
\item Có bao nhiêu cặp $(m,n)$ với $m<n$ sao cho $m+n$ chia hết cho $3$?
\item Có bao nhiêu cặp $(m,n)$ với $m<n$ sao cho $m+n$ chia hết cho $4$?
\item Có bao nhiêu bộ ba $(m,n,p)$ với $m<n<p$ sao cho $m+n+p$ chia hết cho $3$?
\end{enumerate}
\end{probbox}
\probsol{
\textbf{Phân tích đề.} Ở bài 22, ta chia các số theo tính chẵn lẻ (số dư khi chia cho $2$). Ở đây ta chia theo \emph{số dư khi chia
cho $3$} (hoặc cho $4$). Vì $m<n$ (và $m<n<p$) nên mỗi cặp/bộ ba được đếm một lần, không kể thứ tự.
\begin{recallbox}Số dư của một tổng khi chia cho $3$ bằng số dư của tổng các số dư. Ví dụ $7+11=18$: $7$ dư $1$, $11$ dư $2$, $1+2=3$
chia hết cho $3$, nên $18$ chia hết cho $3$.\end{recallbox}

\step{1}{Chia $20$ số theo số dư khi chia cho $3$}
\begin{center}\begin{tabular}{|c|l|c|}\hline
Số dư & Các số & Số lượng\\\hline
$0$ & $3,6,9,12,15,18$ & $6$\\\hline
$1$ & $1,4,7,10,13,16,19$ & $7$\\\hline
$2$ & $2,5,8,11,14,17,20$ & $7$\\\hline
\end{tabular}\end{center}

\step{2}{Câu a}
Tổng hai số dư chia hết cho $3$ chỉ trong hai trường hợp: $0+0$ hoặc $1+2$.
\begin{itemize}[leftmargin=1.2em,itemsep=1pt,topsep=2pt]
\item Hai số cùng dư $0$: chọn $2$ trong $6$ số: $\dfrac{6\times5}{2}=15$ cặp.
\item Một số dư $1$, một số dư $2$: $7\times7=49$ cặp (hai số khác nhóm nên tự động khác nhau; số nhỏ hơn là $m$).
\end{itemize}
(Hai số cùng dư $1$ có tổng dư $2$; hai số cùng dư $2$ có tổng $4$, dư $1$: không thỏa mãn.) Tổng: $15+49=64$ cặp.

\step{3}{Câu b: chia theo số dư khi chia cho $4$}
Mỗi số dư $0,1,2,3$ có đúng $5$ số (vì $20=4\times5$). Tổng hai số dư chia hết cho $4$ khi: $0+0$, $2+2$ hoặc $1+3$.
\[\frac{5\times4}{2}+\frac{5\times4}{2}+5\times5=10+10+25=45 \text{ cặp.}\]

\step{4}{Câu c: các kiểu số dư của bộ ba}
Liệt kê mọi kiểu số dư (không kể thứ tự) của ba số và tổng của chúng:
\begin{center}\begin{tabular}{|c|c|c|c|c|c|c|c|c|c|c|}\hline
Kiểu & $000$ & $111$ & $222$ & $012$ & $001$ & $002$ & $011$ & $022$ & $112$ & $122$\\\hline
Tổng & $0$ & $3$ & $6$ & $3$ & $1$ & $2$ & $2$ & $4$ & $4$ & $5$\\\hline
Chia hết cho $3$? & có & có & có & có & không & không & không & không & không & không\\\hline
\end{tabular}\end{center}
Đếm bốn kiểu thỏa mãn:
\begin{itemize}[leftmargin=1.2em,itemsep=1pt,topsep=2pt]
\item Kiểu $000$: chọn $3$ trong $6$ số: $\dfrac{6\times5\times4}{3\times2\times1}=20$.
\item Kiểu $111$: chọn $3$ trong $7$ số: $\dfrac{7\times6\times5}{3\times2\times1}=35$. Kiểu $222$: cũng $35$.
\item Kiểu $012$: mỗi nhóm một số: $6\times7\times7=294$.
\end{itemize}
Tổng: $20+35+35+294=384$ bộ ba.
\begin{recallbox}Số cách chọn $3$ vật (không kể thứ tự) từ $k$ vật là $\dfrac{k\times(k-1)\times(k-2)}{6}$: chọn có thứ tự được
$k(k-1)(k-2)$ cách, mỗi bộ ba bị đếm $3\times2\times1=6$ lần.\end{recallbox}
\begin{warnbox}Ở kiểu $012$, ta \emph{không} chia cho $6$: ba số thuộc ba nhóm khác nhau, mỗi cách chọn (một số dư $0$, một số dư $1$,
một số dư $2$) cho đúng một bộ ba, sau đó xếp tăng dần là xong.\end{warnbox}
\verify{Tổng số bộ ba là $\dfrac{20\times19\times18}{6}=1140$. Có $384$ bộ có tổng chia hết cho $3$, gần bằng $\dfrac13$ của $1140$ ($380$),
hợp lý vì tổng ba số ``khá đều'' giữa ba số dư.}
\ans{a) $64$ cặp;\quad b) $45$ cặp;\quad c) $384$ bộ ba}
}

\begin{probbox}{48}{\stars{1}}
Với một số có hai chữ số $n=\nd{ab}$, gọi $\nd{ba}$ là \emph{số đảo} của $n$ (như bài 23).
\begin{enumerate}[label=\alph*),leftmargin=1.6em,itemsep=1pt,topsep=2pt]
\item Tìm tất cả các số có hai chữ số $n$ sao cho $n$ cộng với số đảo của nó là một số chính phương.
\item Tìm tất cả các số có hai chữ số, cả hai chữ số đều khác $0$, sao cho $n$ trừ số đảo của nó là một số chính phương dương.
\item Với số có ba chữ số $\nd{abc}$ ($a\ne0$, $c\ne0$), chứng minh rằng hiệu $\nd{abc}-\nd{cba}$ luôn chia hết cho $99$.
\end{enumerate}
\end{probbox}
\probsol{
\textbf{Phân tích đề.} Viết số theo các chữ số: $\nd{ab}=10a+b$ và $\nd{ba}=10b+a$, với $a$ là chữ số hàng chục ($1\le a\le9$),
$b$ là chữ số hàng đơn vị. Số chính phương là bình phương của một số tự nhiên: $1,4,9,16,25,\dots$

\step{1}{Câu a: tính tổng}
$\nd{ab}+\nd{ba}=(10a+b)+(10b+a)=11a+11b=11(a+b)$.
\begin{proofbox}Nếu $11(a+b)=s\times s$ thì $s\times s$ chia hết cho số nguyên tố $11$, nên $s$ chia hết cho $11$ (một tích hai số
chia hết cho số nguyên tố $p$ thì một trong hai số chia hết cho $p$). Khi đó $s\times s$ chia hết cho $121$, tức $11(a+b)$ chia hết
cho $121$, nên $a+b$ chia hết cho $11$.\end{proofbox}
Đối với số có hai chữ số (đảo có thể bắt đầu bằng $0$, như $\nd{ba}=b$ khi $b=0$) ta có $1\le a+b\le18$, nên $a+b=11$, và khi đó
tổng là $11\times11=121=11^2$, đúng là số chính phương. Các số cần tìm: $a+b=11$ với $a$ từ $2$ đến $9$:
\[29,\ 38,\ 47,\ 56,\ 65,\ 74,\ 83,\ 92 \quad(8 \text{ số}).\]

\step{2}{Câu b: tính hiệu}
$\nd{ab}-\nd{ba}=(10a+b)-(10b+a)=9a-9b=9(a-b)$. Hiệu dương nên $a>b\ge1$.
Vì $9=3\times3$ là số chính phương, $9(a-b)$ là số chính phương khi và chỉ khi $a-b$ là số chính phương.
\begin{proofbox}Nếu $a-b=t\times t$ thì $9(a-b)=(3t)\times(3t)$. Ngược lại, nếu $9(a-b)=s\times s$ thì $s\times s$ chia hết cho $3$, nên
$s$ chia hết cho $3$; viết $s=3t$ ta được $9(a-b)=9t\times t$, tức $a-b=t\times t$.\end{proofbox}
Do $1\le b<a\le9$ nên $1\le a-b\le8$; các số chính phương trong khoảng này là $1$ và $4$.
\begin{itemize}[leftmargin=1.2em,itemsep=1pt,topsep=2pt]
\item $a-b=1$: $21, 32, 43, 54, 65, 76, 87, 98$ ($8$ số).
\item $a-b=4$: $51, 62, 73, 84, 95$ ($5$ số).
\end{itemize}
Tất cả có $13$ số.

\step{3}{Câu c: tách theo hàng}
Viết hai số theo giá trị các hàng: $\nd{abc}=100a+10b+c$ và $\nd{cba}=100c+10b+a$. Khi đó
\[\nd{abc}-\nd{cba}=(100a+10b+c)-(100c+10b+a)=100a-a+10b-10b+c-100c,\]
\[\text{tức là}\quad\nd{abc}-\nd{cba}=99a-99c=99(a-c).\]
Vì $a-c$ là một số nguyên nên $99(a-c)$ chia hết cho $99$. (Chữ số hàng chục $b$ triệt tiêu hoàn toàn!)
\begin{warnbox}$\nd{ab}$ là cách viết số có chữ số $a$ và $b$, \emph{không phải} tích $a\times b$. Luôn đổi $\nd{ab}$ thành $10a+b$ trước khi tính.\end{warnbox}
\verify{a) $47+74=121=11^2$. b) $73-37=36=6^2$, $98-89=9=3^2$. c) $731-137=594=99\times6$; $205-502=-297=99\times(-3)$.}
\ans{a) $29$, $38$, $47$, $56$, $65$, $74$, $83$, $92$;\quad b) $13$ số;\quad \mbox{c) $\nd{abc}-\nd{cba}=99(a-c)$.}}
}

\needspace{16\baselineskip}
\begin{probbox}{49}{\stars{2}}
Cho một lưới ô vuông gồm $3$ hàng và $5$ cột ($15$ ô vuông nhỏ) như hình vẽ; ô ở chính giữa được tô màu.
\begin{figC}\begin{scope}[scale=0.6]
  \fill[tmOrange!45] (2,1) rectangle (3,2);
  \draw[Ink,thick] (0,0) grid (5,3);
\end{scope}\end{figC}
\begin{enumerate}[label=\alph*),leftmargin=1.6em,itemsep=1pt,topsep=2pt]
\item Trong hình có bao nhiêu hình chữ nhật (kể cả hình vuông) có các cạnh nằm trên các đường kẻ của lưới?
\item Trong số đó có bao nhiêu hình vuông?
\item Chọn ngẫu nhiên một trong các hình chữ nhật ở câu a. Tính xác suất để hình chữ nhật đó chứa ô được tô màu.
\end{enumerate}
\end{probbox}
\probsol{
\textbf{Phân tích đề.} Giống bài 24, mỗi hình chữ nhật được xác định bởi hai đường kẻ ngang (cạnh trên, cạnh dưới) và hai đường kẻ dọc
(cạnh trái, cạnh phải). Lưới $3$ hàng có $4$ đường ngang, đánh số $0,1,2,3$ từ trên xuống; lưới $5$ cột có $6$ đường dọc, đánh số
$0,1,\dots,5$ từ trái sang.

\step{1}{Câu a: đếm hình chữ nhật}
Chọn $2$ trong $4$ đường ngang: $\dfrac{4\times3}{2}=6$ cách. Chọn $2$ trong $6$ đường dọc: $\dfrac{6\times5}{2}=15$ cách.
Theo quy tắc nhân: $6\times15=90$ hình chữ nhật.

\step{2}{Câu b: đếm hình vuông theo kích thước}
\begin{recallbox}Trong một lưới $3\times5$, hình vuông cạnh $k$ có $(3-k+1)$ vị trí theo chiều dọc và $(5-k+1)$ vị trí theo chiều ngang.\end{recallbox}
\begin{itemize}[leftmargin=1.2em,itemsep=1pt,topsep=2pt]
\item cạnh $1$: $3\times5=15$ hình;
\item cạnh $2$: $2\times4=8$ hình;
\item cạnh $3$: $1\times3=3$ hình;
\item cạnh $4$ trở lên: không có (lưới chỉ cao $3$).
\end{itemize}
Tổng: $15+8+3=26$ hình vuông.
\begin{warnbox}Không được dùng công thức ``chọn $2$ đường ngang và $2$ đường dọc'' để đếm hình vuông: phần lớn các cách chọn đó cho hình chữ
nhật không vuông. Hình vuông phải đếm theo cạnh.\end{warnbox}

\step{3}{Câu c: hình chữ nhật chứa ô ở giữa}
Ô ở giữa nằm giữa đường ngang $1$ và $2$, giữa đường dọc $2$ và $3$. Một hình chữ nhật chứa ô này khi và chỉ khi:
\begin{itemize}[leftmargin=1.2em,itemsep=1pt,topsep=2pt]
\item cạnh trên là đường ngang $0$ hoặc $1$ (ở trên ô): $2$ cách; cạnh dưới là đường ngang $2$ hoặc $3$: $2$ cách;
\item cạnh trái là đường dọc $0$, $1$ hoặc $2$: $3$ cách; cạnh phải là đường dọc $3$, $4$ hoặc $5$: $3$ cách.
\end{itemize}
\begin{figC}\begin{scope}[scale=0.6]
  \fill[tmGreen!25] (1,0) rectangle (4,2);
  \fill[tmOrange!55] (2,1) rectangle (3,2);
  \draw[Ink,thick] (0,0) grid (5,3);
  \draw[tmGreen!70!black,line width=2pt] (1,0) rectangle (4,2);
  \foreach \i in {0,...,5}{\node[font=\scriptsize,text=Slate] at (\i,3.35) {\i};}
  \foreach \j/\l in {3/0,2/1,1/2,0/3}{\node[font=\scriptsize,text=Slate] at (-0.35,\j) {\l};}
\end{scope}\end{figC}
(Ví dụ hình tô xanh lá: cạnh trên là đường ngang $1$, cạnh dưới là đường ngang $3$, cạnh trái là đường dọc $1$, cạnh phải là đường dọc $4$.)
Số hình chữ nhật chứa ô giữa: $2\times2\times3\times3=36$.
\[P=\frac{36}{90}=\frac{2}{5}.\]
\begin{proofbox}Vì sao điều kiện trên là đúng? Hình chữ nhật chứa ô giữa khi và chỉ khi nó ``trùm'' ô đó theo cả hai chiều: cạnh trên phải
ở trên (hoặc trùng) cạnh trên của ô, cạnh dưới ở dưới (hoặc trùng) cạnh dưới của ô; tương tự với cạnh trái, cạnh phải. Bốn lựa chọn này
độc lập với nhau nên ta nhân.\end{proofbox}
\verify{Tách theo từng chiều: trong $15$ cách chọn hai đường dọc, có $3\times3=9$ cách ``trùm'' cột giữa; trong $6$ cách chọn hai đường
ngang, có $2\times2=4$ cách trùm hàng giữa. Xác suất $=\dfrac{9}{15}\times\dfrac{4}{6}=\dfrac35\times\dfrac23=\dfrac25$. Khớp.}
\ans{a) $90$;\quad b) $26$;\quad c) $\dfrac25$}
}

\begin{probbox}{50}{\stars{2}}
Trong một cái hộp có $2026$ viên bi đỏ và $2027$ viên bi xanh. Bên ngoài hộp có một đống bi xanh rất lớn. Bạn Felix lặp lại thao tác
sau: lấy ngẫu nhiên $2$ viên bi trong hộp ra.
\begin{itemize}[leftmargin=1.2em,itemsep=1pt,topsep=2pt]
\item Nếu hai viên \emph{cùng màu} (cùng đỏ hoặc cùng xanh): bỏ cả hai viên ra ngoài, rồi lấy $1$ viên bi xanh từ đống bên ngoài bỏ vào hộp.
\item Nếu hai viên \emph{khác màu}: bỏ viên xanh ra ngoài và trả viên đỏ vào hộp.
\end{itemize}
Felix dừng lại khi trong hộp chỉ còn $1$ viên bi.
\begin{enumerate}[label=\alph*),leftmargin=1.6em,itemsep=1pt,topsep=2pt]
\item Felix đã thực hiện bao nhiêu lần thao tác?
\item Tính xác suất để viên bi cuối cùng là bi đỏ.
\item Nếu lúc đầu hộp có $2025$ viên bi đỏ và $2027$ viên bi xanh thì xác suất ở câu b bằng bao nhiêu?
\end{enumerate}
\end{probbox}
\probsol{
\textbf{Phân tích đề.} Có hàng nghìn viên bi, không thể vẽ cây xác suất như bài 25. Ta tìm một đại lượng \emph{không thay đổi} (gọi là
\emph{bất biến}) trong mỗi thao tác. Ở bài 25, bất biến là ``luôn còn bi xanh''; ở đây ta xét \emph{tính chẵn lẻ của số bi đỏ}.

\step{1}{Câu a: mỗi thao tác làm giảm số bi đi đúng $1$}
\begin{itemize}[leftmargin=1.2em,itemsep=1pt,topsep=2pt]
\item Cùng màu: lấy ra $2$, bỏ vào $1$ $\Rightarrow$ số bi trong hộp giảm $1$.
\item Khác màu: lấy ra $2$, trả lại $1$ $\Rightarrow$ số bi giảm $1$.
\end{itemize}
Lúc đầu có $2026+2027=4053$ viên, cuối cùng còn $1$ viên. Số thao tác là $4053-1=4052$.

\step{2}{Theo dõi số bi đỏ trong mỗi trường hợp}
\begin{center}\begin{tabular}{|c|c|c|}\hline
Lấy được & Số bi đỏ thay đổi & Số bi xanh thay đổi\\\hline
đỏ, đỏ & giảm $2$ & tăng $1$\\\hline
xanh, xanh & không đổi & giảm $2$, tăng $1$: giảm $1$\\\hline
đỏ, xanh & không đổi (trả viên đỏ lại) & giảm $1$\\\hline
\end{tabular}\end{center}
\begin{proofbox}Số bi đỏ hoặc giữ nguyên, hoặc giảm đúng $2$. Bớt đi $2$ không làm thay đổi tính chẵn lẻ (số chẵn trừ $2$ vẫn chẵn,
số lẻ trừ $2$ vẫn lẻ). Vậy \textbf{tính chẵn lẻ của số bi đỏ không bao giờ thay đổi}, dù Felix lấy được những viên nào.\end{proofbox}

\step{3}{Câu b: lúc đầu có $2026$ bi đỏ (số chẵn)}
Số bi đỏ luôn là số chẵn. Khi hộp chỉ còn $1$ viên: nếu viên đó đỏ thì số bi đỏ là $1$ (số lẻ) -- vô lí. Vậy viên cuối cùng chắc chắn
là bi xanh. Biến cố ``viên cuối là bi đỏ'' là biến cố không thể:
\[P(\text{viên cuối đỏ})=0.\]

\step{4}{Câu c: lúc đầu có $2025$ bi đỏ (số lẻ)}
Số bi đỏ luôn là số lẻ, nên không bao giờ bằng $0$. Khi còn $1$ viên, số bi đỏ là số lẻ không vượt quá $1$, tức bằng $1$: viên cuối chắc
chắn là bi đỏ. Biến cố này là biến cố chắc chắn:
\[P(\text{viên cuối đỏ})=1.\]
\begin{tipbox}Số bi xanh ban đầu không ảnh hưởng gì tới kết quả! Chỉ tính chẵn lẻ của số bi đỏ quyết định màu viên bi cuối cùng.\end{tipbox}
\begin{parentbox}Hãy cùng con chơi thử với ít bi (ví dụ $2$ đỏ, $3$ xanh, hoặc $3$ đỏ, $2$ xanh) bằng các đồng xu hai màu. Lần nào cũng thấy
kết quả giống nhau, con sẽ tò mò ``vì sao''; khi đó hãy gợi ý: ``Con thử đếm số bi đỏ sau mỗi lượt xem nó chẵn hay lẻ?''\end{parentbox}
\begin{warnbox}Thấy chữ ``ngẫu nhiên'' rồi nghĩ đáp số phải là một phân số ``đẹp'' như $\frac12$ là sai lầm. Có những quá trình ngẫu nhiên mà kết quả cuối
cùng lại hoàn toàn xác định.\end{warnbox}
\verify{Thử với $2$ đỏ, $1$ xanh: lấy đỏ--đỏ $\to(0\text{ đỏ},2\text{ xanh})\to(0,1)$: xanh; lấy đỏ--xanh $\to(2,0)\to$ lấy đỏ--đỏ $\to(0,1)$: xanh.
Mọi cách đều kết thúc bằng bi xanh, đúng với số bi đỏ ban đầu chẵn.}
\ans{a) $4052$ lần;\quad b) $0$;\quad c) $1$}
}
